\documentclass[11pt]{amsart}
\usepackage{amsmath,amssymb,amsthm}
\usepackage[margin=1.15in]{geometry}
\usepackage[hidelinks]{hyperref}

\newtheorem{theorem}{Theorem}[section]
\newtheorem{lemma}[theorem]{Lemma}
\newtheorem{corollary}[theorem]{Corollary}
\theoremstyle{remark}
\newtheorem{remark}[theorem]{Remark}

\newcommand{\arxiv}[1]{\href{https://arxiv.org/abs/#1}{arXiv:\nolinkurl{#1}}}
\DeclareMathOperator{\conv}{conv}
\setlength{\emergencystretch}{2.5em}

\begin{document}

\title[The Schneider index of Minkowski sums on the line]{The Schneider non-convexity index of Minkowski sums on the
line}
\author{Yitzchak Shmalo}
\date{26 November 2026}
\hypersetup{pdftitle={The Schneider non-convexity index of Minkowski sums on the line},pdfauthor={Yitzchak Shmalo}}
\address{Einstein Institute of Mathematics, The Hebrew University of Jerusalem, Jerusalem, Israel}
\email{yspennstate@gmail.com}
\subjclass[2020]{52A40, 11B13}
\keywords{Schneider non-convexity index, Minkowski sum, sumset, convexification}

\begin{abstract}
The Schneider non-convexity index of a compact set $A\subset\mathbb{R}^n$ is the least $\lambda\ge0$ for which
$A+\lambda\conv A$ is convex. For compact sets $A_1,\dots,A_k\subset\mathbb{R}$ with prescribed indices
$c_1\le\dots\le c_k$ we show that the index of $A_1+\dots+A_k$ is at least $(1-B)_+/(2^k-1+B)$, where
$B=\sum_{i=1}^k 2^{k-i}(1-c_i)/(1+c_i)$. Meyer constructed sets attaining this value and proved the bound for $k=2$;
the bound for every $k$ answers his question on the minimal index of a $k$-fold sum on the line. As a consequence, the
least number of compact subsets of the line with common index $c<1$ whose sum is convex is
$1+\lceil\log_2(1/(1-c))\rceil$, which Meyer derived under the assumption that his sets attain the minimum.
\end{abstract}

\maketitle

\section{Introduction}

For a nonempty compact set $A\subset\mathbb{R}^n$, Schneider~\cite{Sch} defined
\[
  c(A)=\inf\{\lambda\ge0:\ A+\lambda\conv A \text{ is convex}\}.
\]
It satisfies $0\le c(A)\le n$, it vanishes exactly on convex sets, and it is invariant under affine bijections. Its
behaviour under Minkowski addition is part of the study of the convexifying effect of summation, surveyed by
Fradelizi, Madiman, Marsiglietti and Zvavitch~\cite{FMMZ}. Meyer~\cite{MeyerArx,MeyerJIA} asked how small the index
of a sum can be when the indices of the summands are prescribed. For two compact subsets of the line he found the exact
answer~\cite[Theorem~1.2]{MeyerJIA}. For $k$ summands on the line he posed the problem of determining
\[
  F_1(c_1,\dots,c_k)=\min\{c(A_1+\dots+A_k):\ A_i\subset\mathbb{R}\text{ nonempty compact},\ c(A_i)=c_i\}
\]
in~\cite[Problem~2.2]{MeyerArx}, and constructed sets that give an upper bound~\cite[Proposition~2.1]{MeyerArx}. We
prove that his upper bound is the answer.

Throughout, $r(c)=(1-c)/(1+c)$, which maps $[0,1]$ onto $[0,1]$ and is decreasing, and $u_+=\max(u,0)$.

\begin{theorem}\label{thm:main}
Let $k\ge1$ and $0\le c_1\le c_2\le\dots\le c_k\le1$. Put $r_i=r(c_i)$, $M=2^k-1$ and
\[
  B=\sum_{i=1}^k 2^{k-i}r_i .
\]
If $A_1,\dots,A_k\subset\mathbb{R}$ are nonempty compact sets with $c(A_i)=c_i$, then
\[
  c(A_1+\dots+A_k)\ \ge\ \frac{(1-B)_+}{M+B},
\]
and equality holds for the sets of Meyer described in Section~\ref{sec:attain}. Hence
$F_1(c_1,\dots,c_k)=(1-B)_+/(M+B)$.
\end{theorem}

Since $r$ is decreasing, $r_1\ge\dots\ge r_k$, so $B$ is the largest value of $\sum_j2^{j-1}r_{\sigma(j)}$ over
permutations $\sigma$, and the theorem says that the inequality of \cite[Proposition~2.1]{MeyerArx} is an equality.
For $k=2$ it is \cite[Theorem~1.2]{MeyerJIA}. For example, prescribed indices $(3/4,1,1)$ give $B=4/7$ and minimal
index $3/53$, while $(1/2,1,1)$ give $B=4/3$, so a sum of three such sets can be an interval. For equal indices the theorem reads
$F_1(c,\dots,c)=\bigl(1-2^{k-1}(1-c)\bigr)_+/(2^k-1)$.

Meyer also introduced the induced measure of non-convexity: for $x\in[0,n]$ let $N_x$ be the least $k$ such that there
are compact sets $A_1,\dots,A_k\subset\mathbb{R}^n$, each of index $x$, whose sum is convex, and set
$c^\ast(A)=N_{c(A)}-1$~\cite[Definition~2.1]{MeyerArx}. He showed that if equality holds in his Proposition~2.1, then
$c^\ast(A)=\lceil\log_2(1/(1-c(A)))\rceil$ for compact $A\subset\mathbb{R}$ with $c(A)<1$
\cite[Proposition~2.2]{MeyerArx}. Theorem~\ref{thm:main} removes the assumption.

\begin{corollary}\label{cor:induced}
For $0\le x<1$ and $n=1$, $N_x=1+\lceil\log_2(1/(1-x))\rceil$, and $N_1=\infty$. Consequently every nonempty compact
$A\subset\mathbb{R}$ with $c(A)<1$ has $c^\ast(A)=\lceil\log_2(1/(1-c(A)))\rceil$, and $c^\ast(A)=\infty$ if $c(A)=1$.
\end{corollary}

\section{Gaps and the index on the line}

Let $A\subset\mathbb{R}$ be nonempty and compact with $\conv A=[a,b]$, and write $D(A)=b-a$. A \emph{gap} of $A$ is a
bounded component of $\mathbb{R}\setminus A$, that is, an open interval $(u,v)$ with $u,v\in A$ and $(u,v)\cap A=\emptyset$.
Let $g(A)$ be the supremum of the lengths of the gaps of $A$, with $g(A)=0$ if there are none.

\begin{lemma}\label{lem:gap}
Let $h\ge0$. Then $A+[0,h]=[a,b+h]$ if every gap of $A$ has length at most $h$, and $A+[0,h]$ is not convex otherwise.
If $g(A)>0$, some gap has length $g(A)$. Consequently $c(A)=g(A)/D(A)$ when $D(A)>0$, and $c(A)=0$ when $A$ is a point.
\end{lemma}

\begin{proof}
Suppose every gap has length at most $h$, and let $x\in[a,b+h]$. The point $y=\max(A\cap[a,x])$ exists. If $y=x$ or
$y=b$ then $x\in A+[0,h]$. Otherwise $y<x<b$, so $y<\min(A\cap(x,b])=:y'$ and $(y,y')$ is a gap containing $x$; then
$x-y<y'-y\le h$. Conversely, if $(u,v)$ is a gap with $v-u>h$, the point $x=u+h+\tfrac12(v-u-h)$ lies in
$[a,b+h]$, and any $y\in A$ with $x-h\le y\le x$ would lie in $(u,v)$; so $A+[0,h]$ misses $x$ while it contains $a$
and $b+h$.

The gaps are disjoint subintervals of $[a,b]$, so only finitely many have length greater than
$g(A)/2$, and the supremum over those is a maximum. Finally $A+\lambda\conv A$ is a translate of $A+[0,\lambda D(A)]$,
so it is convex exactly when $\lambda D(A)\ge g(A)$.
\end{proof}

\section{The lower bound}

Let $A_1,\dots,A_k$ be as in Theorem~\ref{thm:main}. The index is invariant under translation, so we may assume
$\conv A_i=[0,D_i]$. Put $S=A_1+\dots+A_k$, so that $\conv S=[0,D]$ with $D=\sum_iD_i$, and let $G=g(S)$ and $T=D+G$.
If $D=0$ every $A_i$ is a point, $c_i=0$ for all $i$, $B=M\ge1$ and both sides of the inequality are $0$. Assume $D>0$.

For each $i$ choose $0\le p_i\le q_i\le D_i$ with $q_i-p_i=g(A_i)$ and $(p_i,q_i)\cap A_i=\emptyset$: a gap of
maximal length when $g(A_i)>0$ (Lemma~\ref{lem:gap}), and $p_i=q_i=0$ otherwise. Put $L_i=[0,p_i]$ and
$R_i=[q_i,D_i]$. Then $A_i\subset L_i\cup R_i$, and since $g(A_i)=c_iD_i$ by Lemma~\ref{lem:gap} (also when $D_i=0$),
\begin{equation}\label{eq:sides}
  |L_i|+|R_i|=(1-c_i)D_i .
\end{equation}

\begin{lemma}\label{lem:cover}
Let $w_i=\tfrac12(1+c_i)D_i$. For every nonempty $J\subset\{1,\dots,k\}$,
\[
  \sum_{i\in J}w_i\ \le\ \bigl(1-2^{-|J|}\bigr)\,T .
\]
\end{lemma}

\begin{proof}
By Lemma~\ref{lem:gap}, $S+[0,G]=[0,T]$. For each choice $\varepsilon=(\varepsilon_i)_{i\in J}$ with
$\varepsilon_i\in\{L_i,R_i\}$ let
\[
  I_\varepsilon=\sum_{i\in J}\varepsilon_i+\sum_{i\notin J}[0,D_i]+[0,G],
\]
a closed interval of length $\sum_{i\in J}|\varepsilon_i|+\sum_{i\notin J}D_i+G$. Every point of $S$ has the form
$\sum_ia_i$ with $a_i\in A_i\subset L_i\cup R_i$, so $S+[0,G]\subset\bigcup_\varepsilon I_\varepsilon$. There are
$2^{|J|}$ choices, and for each $i\in J$ each of $L_i,R_i$ occurs in exactly half of them. By~\eqref{eq:sides},
\[
  T\le\sum_\varepsilon|I_\varepsilon|
   =2^{|J|}\Bigl(G+\sum_{i\notin J}D_i\Bigr)+2^{|J|-1}\sum_{i\in J}(1-c_i)D_i
   =2^{|J|}\Bigl(T-\sum_{i\in J}w_i\Bigr),
\]
where the last step uses $G+\sum_{i\notin J}D_i=T-\sum_{i\in J}D_i$.
\end{proof}

\begin{proof}[Proof of the inequality in Theorem~\ref{thm:main}]
Put $v_i=2/(1+c_i)=1+r_i$ for $1\le i\le k$ and $v_{k+1}=0$. Since the $c_i$ increase, $v_1\ge v_2\ge\dots\ge
v_{k+1}$. With $W_j=\sum_{i\le j}w_i$, summation by parts and Lemma~\ref{lem:cover} for the sets $\{1,\dots,j\}$ give
\[
  D=\sum_{i=1}^kv_iw_i=\sum_{j=1}^k(v_j-v_{j+1})W_j
   \le T\sum_{j=1}^k(v_j-v_{j+1})\bigl(1-2^{-j}\bigr)
   =T\sum_{i=1}^k2^{-i}v_i .
\]
Here $\sum_i2^{-i}v_i=\sum_i2^{-i}+\sum_i2^{-i}r_i=2^{-k}(M+B)$. Hence $2^kD\le(D+G)(M+B)$, that is,
\[
  \frac GD\ \ge\ \frac{2^k-M-B}{M+B}=\frac{1-B}{M+B}.
\]
By Lemma~\ref{lem:gap}, $c(S)=G/D$, and $c(S)\ge0$.
\end{proof}

\begin{remark}
Only the $k$ initial segments $\{1,\dots,j\}$ of the order of the indices enter the proof; the order of the diameters
$D_i$ plays no role.
\end{remark}

\section{Attainment and the induced measure}\label{sec:attain}

The following sets are those of \cite[Proposition~2.1]{MeyerArx}, listed in the order of the indices. Let
$h_i=2^{k-i}$ and
\[
  A_i=[0,h_ir_i]\cup[h_i,h_i(1+r_i)]=\{0,h_i\}+[0,h_ir_i].
\]
If $c_i<1$ then $r_i>0$, $D(A_i)=h_i(1+r_i)$ and $g(A_i)=h_i(1-r_i)$, and if $c_i=1$ then $A_i=\{0,h_i\}$; in both
cases $c(A_i)=(1-r_i)/(1+r_i)=c_i$ by Lemma~\ref{lem:gap}. The subset sums of $h_1,\dots,h_k$ are exactly
$0,1,\dots,M$, so
\[
  A_1+\dots+A_k=\{0,1,\dots,M\}+[0,B]=\bigcup_{j=0}^M[j,j+B].
\]
This set has diameter $M+B$, and all its gaps have length $1-B$ if $B<1$, while it is the interval $[0,M+B]$ if $B\ge1$.
Its index is $(1-B)_+/(M+B)$, which completes the proof of Theorem~\ref{thm:main}.

\begin{proof}[Proof of Corollary~\ref{cor:induced}]
Take $c_1=\dots=c_k=x$ and $r=r(x)$, so $B=Mr$. A sum of $k$ compact sets of index $x$ can be convex if and only if
$F_1(x,\dots,x)=0$, because the minimum in Theorem~\ref{thm:main} is attained. If $x=1$ then $B=0$ and the minimum is
$1/M>0$ for every $k$, so $N_1=\infty$. If $x<1$, the minimum vanishes if and only if $(2^k-1)(1-x)\ge1+x$, which
rearranges to $2^{k-1}\ge1/(1-x)$, that is, $k\ge1+\log_2(1/(1-x))$. The least such integer is
$1+\lceil\log_2(1/(1-x))\rceil$.
\end{proof}

\section*{Acknowledgements and funding}
This work was done by AI agents, Claude (Anthropic) and Codex (OpenAI), connected to Math Brain~\cite{MathBrainRepo}
and working together inside the Brain framework~\cite{BrainFramework}, a shared, persistent memory in which the agents
record claims, proofs, failed attempts and methods, retrieve one another's results, coordinate who works on what, and
audit each other's work. For more on the Brain and Math Brain, see the Math Brain paper~\cite{MathBrainPaper}, the Math
Brain repository~\cite{MathBrainRepo} and the Brain framework repository~\cite{BrainFramework}.

The research presented in this paper was supported by the European Research Council (ERC) under the European Union's
Horizon Europe research and innovation programme (grant agreement No.~101041711), by the Simons Foundation as part of
the Collaboration on the Mathematical and Scientific Foundations of Deep Learning, by Heights Labs, by Convex Nexus
Capital, by the Israel Science Foundation (grant number 2258/19), and by the Israel Science Foundation (ISF Grant
4101/25).

\section*{Declaration of competing interest}
The author holds equity in Heights Labs and in Convex Nexus Capital and has a management role at Convex Nexus Capital.
These entities are included in the funding declaration above.

\section*{Data availability}
No data were generated or analysed in this work. A Python script that checks the inequality of
Theorem~\ref{thm:main} in exact rational arithmetic on random finite unions of intervals with $k=2,3,4$, and checks
that Meyer's sets attain it on a grid of index vectors, is \texttt{code/check\_line\_sums.py} at
\url{https://github.com/yspennstate/schneider-index-line-sums/tree/v1.0}.

\section*{Declaration of generative AI and AI-assisted technologies}
This paper is part of a research project on AI proving problems. Codex agents (OpenAI) produced the proof, and Claude
agents (Anthropic) checked the proof and the references and wrote the manuscript. The problem, the candidate formula
and the extremal sets are due to Meyer~\cite{MeyerArx}. The author directed the work and is responsible for the content
of the paper and for decisions concerning its dissemination.

\begin{thebibliography}{9}
\bibitem{FMMZ} M.~Fradelizi, M.~Madiman, A.~Marsiglietti and A.~Zvavitch, \emph{The convexification effect of
Minkowski summation}, EMS Surv. Math. Sci. \textbf{5} (2018), no.~1--2, 1--64,
\href{https://doi.org/10.4171/EMSS/26}{doi:10.4171/EMSS/26}; \arxiv{1704.05486}.
\bibitem{MeyerArx} M.~Meyer, \emph{Measuring the convexity of compact sumsets with the Schneider non-convexity index},
preprint, 2024, \arxiv{2405.00221v1}.
\bibitem{MeyerJIA} M.~Meyer, \emph{Measuring the convexity of compact sumsets with the Schneider non-convexity index},
J. Inequal. Appl. \textbf{2025} (2025), Paper No.~103,
\href{https://doi.org/10.1186/s13660-025-03356-w}{doi:10.1186/s13660-025-03356-w}.
\bibitem{Sch} R.~Schneider, \emph{A measure of convexity for compact sets}, Pacific J. Math. \textbf{58} (1975),
no.~2, 617--625.
\bibitem{MathBrainRepo} Y.~Shmalo, \emph{Math Brain: public snapshot, code and tests}, 2026,
\url{https://github.com/yspennstate/math-brain-public/tree/5427734deee9be643d6dfbdd5e7d4baeeefe4a4a}.
\bibitem{MathBrainPaper} Y.~Shmalo, \emph{Math Brain: shared, auditable memory for multi-agent mathematical
research}, 2026, \href{https://doi.org/10.5281/zenodo.24150837}{doi:10.5281/zenodo.24150837},
\url{https://github.com/yspennstate/math-brain-public/tree/5427734deee9be643d6dfbdd5e7d4baeeefe4a4a/papers/math-brain-2026-10-24}.
\bibitem{BrainFramework} Y.~Shmalo, \emph{Brain Framework}, public software repository, 2026,
\url{https://github.com/yspennstate/brain-framework-public/tree/46b935d26d5a2893f7790b60466387dd8da47ef5}.
\end{thebibliography}

\end{document}
